\section{Perron dynamics and the stable memory mode}

On a closed route, the cylindrical Perron cocycle satisfies

\begin{equation}
V_n=\varphi^n.
\label{rm:eq:cosmo-perron-V}
\end{equation}

In spatial dimension \(D=3\), the scale determined by the cocycle is

\begin{equation}
\frac{a_n}{a_0}=V_n^{1/3}=\varphi^{n/3},
\label{rm:eq:cosmo-perron-a}
\end{equation}

whereas the stable mode of the quadratic action is

\begin{equation}
M_n=\varphi^{-2n}=V_n^{-2}.
\label{rm:eq:cosmo-perron-M}
\end{equation}

Eliminating \(n\) between \cref{rm:eq:cosmo-perron-a,rm:eq:cosmo-perron-M} gives the
central result:

\begin{theorem}[Law cylindrical of sixth order for the memory]
For every closed level of the Perron realization in \(D=3\),
\begin{equation}
\boxed{
M_n=\left(\frac{a_n}{a_0}\right)^{-6}.}
\label{rm:eq:cosmo-minus-six}
\end{equation}
The exponent \(-6\) is not fitted: it is the quotient of the stable exponent
\(-2\) by the dimensional root \(1/3\).
\end{theorem}

In the distinguished returns one obtains

\[
\frac{a_9}{a_0}=\varphi^3,
\qquad
\frac{a_{27}}{a_0}=\varphi^9,
\qquad
\frac{a_{108}}{a_0}=\varphi^{36},
\]

and the same levels satisfy exactly
\(M_n=(a_n/a_0)^{-6}\). If the additive temporal parameter \(t_n=nt_0\) is declared, the
logarithmic rate is constant:

\begin{equation}
H_\varphi
:=\frac{\Delta\log a_n}{\Delta t}
=\frac{\log\varphi}{3t_0},
\qquad
H_\varphi t_0
=0.1604039416865344824992529711\ldots.
\label{rm:eq:cosmo-Hphi}
\end{equation}

The symbol \(H_\varphi\) names the rate of this discrete scale. Reading it
as a Hubble parameter requires the cosmographic transducer; the internal
identity does not depend on that promotion.

\section{Realization in the system of Einstein--Cartan}

In an Einstein--Cartan spin fluid, the quadratic spin density
obeys

\begin{equation}
s^2\propto a^{-6}.
\label{rm:eq:cosmo-EC-scaling}
\end{equation}

The \cref{rm:eq:cosmo-minus-six} provides the exponent and the stable mode without
consulting a bounce. To make the physical reading, one must declare a
positive section \(s_*^2\) and a map

\begin{equation}
\mathcal T_{\rm EC}:M_n\longmapsto
s_n^2:=s_*^2M_n
=s_*^2\left(\frac{a_n}{a_0}\right)^{-6}.
\label{rm:eq:cosmo-EC-transducer}
\end{equation}

The map preserves the scaling law, but does not by itself determine the amplitude
\(s_*^2\), the spin current, the sign of the term of contorsion, or the
equation of state. With mass density, an effective equation has the
form

\begin{equation}
H^2=\frac{8\pi G}{3}
\bigl(\rho_{\rm m}-\rho_{\rm spin,m}\bigr)+\cdots,
\label{rm:eq:cosmo-EC-mass}
\end{equation}

whereas with energy density it must be written

\begin{equation}
H^2=\frac{8\pi G}{3c^2}
\bigl(\rho_{\rm E}-\rho_{\rm spin,E}\bigr)+\cdots.
\label{rm:eq:cosmo-EC-energy}
\end{equation}

The two conventions are not mixed. A bounce candidate requires, in addition to
equality of the densities, regularity, matter, and boundary data.

\begin{falsifier}
The Einstein--Cartan transduction fails if the scale functional produces an
exponent other than \(-6\), if \(V_n\ne\varphi^n\) along the declared route, or
if the physical map does not preserve positivity and dimensions. Retrospectively choosing
\(s_*^2\) to force a bounce threshold does not establish the
transducer.
\end{falsifier}

\section{Structural comparison between the 108-state cycle and Perron dynamics}

\begin{center}
\begin{tabularx}{\textwidth}{>{\bfseries}p{29mm}XX}
\toprule
& CT108--de Sitter framework & Perron--Einstein--Cartan framework\\
\midrule
Antecedent & circular return of order \(108\) & cylindrical cocycle and stable mode\\
Selection mechanism & \(r_g=R_{108}\) & \(V_n=\varphi^n\), \(D=3\)\\
Magnitudes determinadas & radius, curvature, temperature, and entropy & expansion, rate, and dilution of memory\\
Invariant & \(G_a\hbar_a\) and circular radius & \(M_na_n^6\)\\
Status & de Sitter conditioned interface & exact internal law; conditioned EC reading\\
Falsification criterion & incompatible radius or background equation & perdida of the exponente \(-6\) or of the transductor\\
\bottomrule
\end{tabularx}
\end{center}

\begin{statusbox}
The CT108 identities, the selector \(\theta_{108}\), and the law
\(M_n=(a_n/a_0)^{-6}\) are \status{Internal exact theorem}. The reading of the same radius
as de Sitter and the identification of \(M_n\) with a spin density are
\status{Conditional physical realizations}, each with its own falsification criterion.
\end{statusbox}
